\documentclass[10pt,a4paper]{amsart}
\usepackage[margin=19mm]{geometry}
\usepackage{amsmath,amssymb,mathtools}
\usepackage[scr=rsfs,cal=euler]{mathalfa}
\usepackage{xfrac}
\usepackage[unicode,hidelinks,bookmarks=false]{hyperref}
\newcommand{\ie}{i.e.\ }
\newcommand{\cf}{cf.\ }
\newcommand{\resp}{respectively\ }
\newcommand{\Subsection}{\subsection}
\providecommand{\nobreakdash}{\nobreak\hbox{-}}
\setlength{\emergencystretch}{2em}
\multlinegap0pt
\usepackage{bm}
\usepackage{bbm}
\usepackage{dsfont}
\usepackage{xspace}
\usepackage{cancel}
\usepackage{mathtools}
\usepackage{amsthm}
\newtheorem{theorem}{Theorem}
\newtheorem{fact}[theorem]{Fact}
\newcommand{\perf}{\infty}
\title[Remarks on some Homological Problems regarding Infinite Inte]{Remarks on some Homological Problems regarding Infinite Integral Extensions}
\author{Mohsen Asgharzadeh, Shravan Patankar}
\date{Source: arXiv:2608.09473v2 (CC BY 4.0)}
\begin{document}
\maketitle

\noindent\emph{Source and licence.} Remarks on some Homological Problems regarding Infinite Integral Extensions. Version arXiv:2608.09473v2 (CC BY 4.0), distributed under the Creative Commons Attribution 4.0 International licence, as recorded in the arXiv licence metadata for this exact version and shown on the abstract page https://arxiv.org/abs/2608.09473v2. Full rights notice: licenses/arxiv-2608.09473-CC-BY-4.0.txt.\par\smallskip

\noindent\textit{Editorial scope.} Counted unit: the Fact whose source label is \texttt{None} in the article (source0.tex lines 977--983, source label \texttt{None}), delivered together with the article's own complete original proof of it (source0.tex lines 985--1057, one \texttt{proof} environment of 73 lines). No mathematics is written, repaired or re-derived: statement and proof are the article's own text line for line. Required context retained uncounted: none. Every definition, display and label of those lines is retained unchanged. Omitted from this package and not redistributed: 1--976, 984--984, 1058--1629. Nothing inside a retained excerpt is omitted or altered: every retained body is delivered exactly as the article's lines are written, so no omission receipt of any kind accompanies this package. Citations: the retained excerpts carry no citation command at all, so no bibliography entry of the article is transcribed and no reference is redistributed. Reference adaptation: none was needed and none was made; every reference inside the delivered unit resolves inside it, so no unresolved reference or citation is produced. Printed slips kept literal: no printed word, symbol, hypothesis or number of the counted statement or of its proof was altered. Layout and class: the article's own class is \texttt{amsart} with packages including geometry, amssymb, stmaryrd, microtype, mathtools, enumitem, cite, hyphenat, color, tikz-cd, hyperref, bookmark, breakurl; the wrapper is the lane's pinned \texttt{amsart} editorial wrapper at 10pt on A4, which restates the theorem-like environments the retained text needs and adds the article's own macro definitions that the retained text needs (\texttt{\textbackslash perf}), with extra wrapper packages (bm, bbm, dsfont, xspace, cancel, mathtools). The delivered numbering is the unit's own. Assets: no retained span contains a figure environment, a graphics-inclusion command, a PDF-page inclusion, a table or a TikZ picture; no image file is copied into this package, the package declares no asset dependency and no image file is redistributed with it. Licence: version arXiv:2608.09473v2 of this article is distributed under the Creative Commons Attribution 4.0 International licence, as recorded in the arXiv licence metadata for this exact version and shown on the abstract page https://arxiv.org/abs/2608.09473v2. A full rights notice is delivered at licenses/arxiv-2608.09473-CC-BY-4.0.txt. Characters: every retained excerpt is pure ASCII, so the delivered engine, XeLaTeX with its default Latin Modern text font, reports no missing character.
	\begin{fact}
		Let \(R\) be a \(p\)-torsion-free, \(p\)-adically complete ring equipped with a \(p\)-derivation$
		\delta:R\to R.$ Put $A=R/pR$ and let $A^{\perf}$ be the perfection of \(A\). Denote by
		$\iota:A\to A^{\perf}$ the canonical map. Then the coperfection of \(R\) in the category of \(p\)-torsion-free, \(p\)-adically complete \(\delta\)-rings is identified,   with
		$W(A^{\perf}).$ Consequently there is a morphism of \(p\)-complete \(\delta\)-rings
		$\theta_\delta:R\to W(A^{\perf}). $ The morphism is an injection if $A$ is reduced.
	\end{fact}

	\begin{proof}
		We use the Joyal, or \(y\)-, coordinates on \(W(A^{\perf})\). Thus an element of \(W(A^{\perf})\) is written
		$
		(y_0,y_1,y_2,\ldots),
		$
		with the property that the \(p\)-derivation on \(W(A^{\perf})\) is given by the shift
		$
		\delta(y_0,y_1,y_2,\ldots)
		=
		(y_1,y_2,y_3,\ldots).
		$
		With this convention, the morphism \(\theta_\delta\) is given by
		$$\theta_\delta(x)
		=
		\left(
		\iota(\overline{x}),
		\iota(\overline{\delta(x)}),
		\iota(\overline{\delta^2(x)}),
		\iota(\overline{\delta^3(x)}),
		\ldots
		\right),
		$$ 
		where \(\overline{x}\) denotes the image of \(x\) in \(R/pR\). Equivalently,
		$
		y_n\bigl(\theta_\delta(x)\bigr)
		=
		\iota (\overline{\delta^n(x)})
		.
		$
		The \(p\)-derivation \(\delta\) determines the Frobenius lift	
		$
		\varphi(x)=x^p+p\delta(x).
		$
		The reduction map
		$
		R\to A,
		$
		followed by \(\iota:A\to A^{\perf}\), gives a ring homomorphism
		$
		R\to A^{\perf},
		.
		$
		Since \(A^{\perf}\) is perfect of characteristic \(p\), tilting identifies the corresponding \(p\)-complete perfect \(\delta\)-ring with
		$
		W(A^{\perf}).
		$
		The resulting lift is the morphism
		$
		\theta_\delta:R\to W(A^{\perf}).
		$ It remains to determine its coordinates. Write
		$
		\theta_\delta(x)
		=
		(y_0,y_1,y_2,\ldots)
		$ in Joyal coordinates. Since the \(p\)-derivation on \(W(A^{\perf})\) is the shift,
		$
		y_n
		=
		y_0\circ\delta^n.
		$
		Because \(\theta_\delta\) is a morphism of \(\delta\)-rings,
		$\delta^n(\theta_\delta(x))=
		\theta_\delta(\delta^n(x)).$
		Taking zeroth coordinates gives
		$y_n(\theta_\delta(x))
		=
		y_0(\theta_\delta(\delta^n(x))).
		$ By the definition of the lift, its zeroth coordinate is reduction modulo \(p\), followed by \(\iota\). Hence
		$y_n(\theta_\delta(x))
		=
		\iota(\overline{\delta^n(x)}).$
		This implies the desired claim.
	\end{proof}

\end{document}
